\chapter{RESULTADOS E DISCUSSÕES}
Este capítulo destina-se à apresentação fiel e objetiva dos resultados alcançados com o desenvolvimento do projeto ou estudo, acompanhada obrigatoriamente da sua discussão crítica. A construção deste capítulo exige o cruzamento de duas ações: a descrição precisa dos resultados obtidos (testes de software, métricas de desempenho, avaliação de usabilidade, aplicação de questionários ou estudo de caso) e a interpretação científica desses dados, confrontando-os com o problema de pesquisa, os objetivos propostos, a fundamentação teórica e os trabalhos correlatos.

Sempre que houver dados comparativos, métricas ou sínteses quantificáveis, utilize tabelas ou quadros apropriados (como exemplificado na Tabela \ref{tab:acompanhamento}), inserindo o título na parte superior e a fonte na parte inferior (Normas ABNT/UEMA):

\begin{table}[htb]
\centering
\caption{\label{tab:acompanhamento}Exemplo de tabela de acompanhamento de atividades}
\begin{tabular}{|l|c|c|c|c|c|c|c|}
\hline
\textbf{Atividades / Meses} & \textbf{Jun} & \textbf{Jul} & \textbf{Ago} & \textbf{Set} & \textbf{Out} & \textbf{Nov} & \textbf{Dez} \\ \hline
Levantamento de requisitos & X & & & & & & \\ \hline
Modelagem do sistema & & X & X & & & & \\ \hline
Implementação do protótipo & & & X & X & & & \\ \hline
Testes e validação & & & & X & X & & \\ \hline
Redação da monografia & & & & & X & X & \\ \hline
Apresentação & & & & & & & X \\ \hline
\end{tabular}
\fonte{Elaborado pelo autor (\the\year).}
\end{table}
